% GENERATED FILE. Edit canonical lessons, then run npm run generate:books.
% canonical-source-hash: fnv1a64:541a82d044a6ef2e
% canonical-lessons: KA-C250-sutta, KA-C250-gudisalu, KA-C250-badavane, KA-C250-ance-kaceri, KA-C250-kanive

\chapter{Around, Hut, Residential area, Post office, Valley}
\label{ch:ka-around-hut-residential-area-post-office-valley}
\hlchaptermodality{250}

\begin{chapteropening}
\textbf{By the end of this chapter:} \emph{I can say around, a hut, a residential area, a post office and a valley.}

Complete the last lesson of chapter 250: I can say around, a hut, a residential area, a post office and a valley.
\end{chapteropening}

\section[\texorpdfstring{\kn{ಸುತ್ತ} (sutta)}{ಸುತ್ತ (sutta)}]{\kn{ಸುತ್ತ} (sutta) — around}

\label{lesson:KA-C250-sutta}



\begin{quote}
Before the new one: say the Kannada for opposite, then the Kannada for a district.
\end{quote}

\subsection*{You'll want to know: \kn{ಸುತ್ತ}}
\textbf{\kn{ಸುತ್ತ}} — \emph{sutta} — \textquotedblleft{}around\textquotedblright{}.

\begin{grammarlens}[title={What you've built}]
One more word for this chapter.
\end{grammarlens}

\subsection*{Your turn}
\begin{itemize}
\raggedright
  \item \emph{Say it:} \emph{sutta}
  \item \emph{Say it:} \emph{sutta}, once more
  \item \emph{Say it:} say \emph{sutta}
  \item \emph{Recall:} say the Kannada for opposite, then the Kannada for a district, then say \emph{sutta} again
\end{itemize}

\begin{tcolorbox}[breakable,colback=teal!4,colframe=teal!35!black,title={Before you move on}]
What does \kn{ಸುತ್ತ} mean? (\textquotedblleft{}Around\textquotedblright{}.) Say it once more.
\end{tcolorbox}
\section[\texorpdfstring{\kn{ಗುಡಿಸಲು} (guḍisalu)}{ಗುಡಿಸಲು (guḍisalu)}]{\kn{ಗುಡಿಸಲು} (guḍisalu) — a hut}

\label{lesson:KA-C250-gudisalu}



\begin{quote}
Before the new one: say the Kannada for a district, then the Kannada for around.
\end{quote}

\subsection*{You'll want to know: \kn{ಗುಡಿಸಲು}}
\textbf{\kn{ಗುಡಿಸಲು}} — \emph{guḍisalu} — \textquotedblleft{}a hut\textquotedblright{}.

\begin{grammarlens}[title={What you've built}]
One more word for this chapter.
\end{grammarlens}

\subsection*{Your turn}
\begin{itemize}
\raggedright
  \item \emph{Say it:} \emph{guḍisalu}
  \item \emph{Say it:} \emph{guḍisalu}, once more
  \item \emph{Say it:} say \emph{guḍisalu}
  \item \emph{Recall:} say the Kannada for a district, then the Kannada for around, then say \emph{guḍisalu} again
\end{itemize}

\begin{tcolorbox}[breakable,colback=teal!4,colframe=teal!35!black,title={Before you move on}]
What does \kn{ಗುಡಿಸಲು} mean? (\textquotedblleft{}A hut\textquotedblright{}.) Say it once more.
\end{tcolorbox}
\section[\texorpdfstring{\kn{ಬಡಾವಣೆ} (baḍāvaṇe)}{ಬಡಾವಣೆ (baḍāvaṇe)}]{\kn{ಬಡಾವಣೆ} (baḍāvaṇe) — a residential area, a neighbourhood}

\label{lesson:KA-C250-badavane}



\begin{quote}
Before the new one: say the Kannada for around, then the Kannada for a hut.
\end{quote}

\subsection*{You'll want to know: \kn{ಬಡಾವಣೆ}}
\textbf{\kn{ಬಡಾವಣೆ}} — \emph{baḍāvaṇe} — \textquotedblleft{}a residential area, a neighbourhood\textquotedblright{}.

\begin{grammarlens}[title={What you've built}]
One more word for this chapter.
\end{grammarlens}

\subsection*{Your turn}
\begin{itemize}
\raggedright
  \item \emph{Say it:} \emph{baḍāvaṇe}
  \item \emph{Say it:} \emph{baḍāvaṇe}, once more
  \item \emph{Say it:} say \emph{baḍāvaṇe}
  \item \emph{Recall:} say the Kannada for around, then the Kannada for a hut, then say \emph{baḍāvaṇe} again
\end{itemize}

\begin{tcolorbox}[breakable,colback=teal!4,colframe=teal!35!black,title={Before you move on}]
What does \kn{ಬಡಾವಣೆ} mean? (\textquotedblleft{}A residential area, a neighbourhood\textquotedblright{}.) Say it once more.
\end{tcolorbox}
\section[\texorpdfstring{\kn{ಅಂಚೆ} \kn{ಕಚೇರಿ} (añce kacēri)}{ಅಂಚೆ ಕಚೇರಿ (añce kacēri)}]{\kn{ಅಂಚೆ} \kn{ಕಚೇರಿ} (añce kacēri) — a post office}

\label{lesson:KA-C250-ance-kaceri}



\begin{quote}
Before the new one: say the Kannada for a hut, then the Kannada for a residential area.
\end{quote}

\subsection*{You'll want to know: \kn{ಅಂಚೆ} \kn{ಕಚೇರಿ}}
\textbf{\kn{ಅಂಚೆ} \kn{ಕಚೇರಿ}} — \emph{añce kacēri} — \textquotedblleft{}a post office\textquotedblright{}.

\begin{grammarlens}[title={What you've built}]
One more word for this chapter.
\end{grammarlens}

\subsection*{Your turn}
\begin{itemize}
\raggedright
  \item \emph{Say it:} \emph{añce kacēri}
  \item \emph{Say it:} \emph{añce kacēri}, once more
  \item \emph{Say it:} say \emph{añce kacēri}
  \item \emph{Recall:} say the Kannada for a hut, then the Kannada for a residential area, then say \emph{añce kacēri} again
\end{itemize}

\begin{tcolorbox}[breakable,colback=teal!4,colframe=teal!35!black,title={Before you move on}]
What does \kn{ಅಂಚೆ} \kn{ಕಚೇರಿ} mean? (\textquotedblleft{}A post office\textquotedblright{}.) Say it once more.
\end{tcolorbox}
\section[\texorpdfstring{\kn{ಕಣಿವೆ} (kaṇive)}{ಕಣಿವೆ (kaṇive)}]{\kn{ಕಣಿವೆ} (kaṇive) — a valley}

\label{lesson:KA-C250-kanive}



\begin{quote}
Before the new one: say the Kannada for a residential area, then the Kannada for a post office.
\end{quote}

\subsection*{You'll want to know: \kn{ಕಣಿವೆ}}
\textbf{\kn{ಕಣಿವೆ}} — \emph{kaṇive} — \textquotedblleft{}a valley\textquotedblright{}.

\begin{grammarlens}[title={What you've built}]
One more word for this chapter.
\end{grammarlens}

\subsection*{Your turn}
\begin{itemize}
\raggedright
  \item \emph{Say it:} \emph{kaṇive}
  \item \emph{Say it:} \emph{kaṇive}, once more
  \item \emph{Say it:} say \emph{kaṇive}
  \item \emph{Recall:} say the Kannada for a residential area, then the Kannada for a post office, then say \emph{kaṇive} again
\end{itemize}

\begin{tcolorbox}[breakable,colback=teal!4,colframe=teal!35!black,title={Before you move on}]
What does \kn{ಕಣಿವೆ} mean? (\textquotedblleft{}A valley\textquotedblright{}.) Say it once more.
\end{tcolorbox}
